\paragraph{Theorem F.2.38 (Proved: the complete ordered involution stratum).}
There are exactly 56 simultaneous $\mathrm{GL}_6(\mathbb Z)$
conjugacy classes of ordered tuples $(M_1,M_2,M_0)$ satisfying
the primitive maximal-cusp conditions of F.2 with both finite
linear orders equal to two.  An explicit integral normalizer
returns a canonical representative and a checked unimodular
conjugating matrix for every such input, without any bound on
its entries.  Every representative has free affine lifts of
actual order four at both finite points.  None can give a
simply connected filling, for any choices of free affine
orders or twists.

\begin{proof}
First choose an integral cusp basis, so that
\[
 M_0=I_6+N,\qquad N=\begin{pmatrix}0&I_3\\0&0\end{pmatrix}.
                                                               \tag{F.IV1}
\]
This is effective by Smith reduction.  In fact choose vectors
$\beta_i$ whose images $\alpha_i=N\beta_i$ form a primitive
basis of $\operatorname{im}N=\ker N$.
The index-one hypothesis says that their classes form a basis
of the quotient, so $[\alpha_1,\alpha_2,\alpha_3,
\beta_1,\beta_2,\beta_3]$ is unimodular.

Since $M_2=M_1(I-N)$ and $M_1^2=M_2^2=I$, we have
$M_1NM_1=-N$.  Solving this block equation gives precisely
\[
 M_1=\begin{pmatrix}A&B\\0&-A\end{pmatrix},\qquad
 M_2=\begin{pmatrix}A&B-A\\0&-A\end{pmatrix},\qquad
 A^2=I_3,\quad AB=BA.                                     \tag{F.IV2}
\]
Every integral conjugator preserving (F.IV1) is
$\begin{pmatrix}C&D\\0&C\end{pmatrix}$, with
$C\in\mathrm{GL}_3(\mathbb Z)$ and $D$ arbitrary integral.
Thus one first classifies $A$ integrally; after fixing it,
$B$ is considered modulo $AD+DA$ and modulo conjugation by
the integral centralizer of $A$.

We include the integral involution classification needed here.
For the lattice $L=\mathbb Z^3$, let $L_+$ and $L_-$ be the
primitive kernels of $A-I$ and $A+I$.  They satisfy
$2L\subseteq L_+\oplus L_-\subseteq L$.
The quotient $L/(L_+\oplus L_-)$ is a vector space over
$\mathbb F_2$.  Its projections into $L_+/2L_+$ and
$L_-/2L_-$ are injective: a glue vector with only one
nonzero projection would give a new integral vector in the
corresponding eigensublattice, contradicting its definition.
It is consequently the graph of an isomorphism between two
subspaces of equal dimension $r$.
Unimodular changes of bases in $L_\pm$ put this graph in
coordinate form.  For each glued pair $e_+,e_-$ the vectors
$(e_++e_-)/2$ and $(e_+-e_-)/2$ form a swap block.
The unglued vectors give trivial and sign blocks.  Hence
there is a unique type
\[
 A\cong I_p\oplus(-I_q)\oplus
 \begin{pmatrix}0&1\\1&0\end{pmatrix}^{\oplus r},
 \qquad p+q+2r=3.                                         \tag{F.IV3}
\]
The dimensions of the two rational eigenspaces and the
index $[L:L_+\oplus L_-]=2^r$ prove uniqueness.
Here $r$ is zero or one, so this normalization only requires
integer kernels, one Smith reduction and, in the glued case,
one explicit half-sum pair whose integrality is checked.

For fixed (F.IV3), compute the quotient of the commuting
matrices by the subgroup $\{AD+DA\}$.
On the trivial and sign diagonal blocks the norm is
respectively $2D$ and $-2D$.  The commuting off-diagonal
blocks between the trivial and sign summands are zero.
Every commuting block involving a swap summand is a norm.
For example a trivial-to-swap block is a row $(a,a)$ and
is obtained as $(d,e)+(e,d)$; a sign-to-swap block has the
same calculation with a difference.  For two swap summands,
the commuting block is
$\begin{pmatrix}a&b\\b&a\end{pmatrix}$, and applying the
norm to $\begin{pmatrix}d&e\\f&g\end{pmatrix}$ gives
$\begin{pmatrix}e+f&d+g\\d+g&e+f\end{pmatrix}$.
The reverse mixed blocks have the transposed calculations.
Therefore the full integral quotient is exactly
\[
 \operatorname{End}_{\mathbb Z[A]}(L)/\{AD+DA\}
 \cong \operatorname{Mat}_p(\mathbb F_2)
       \oplus\operatorname{Mat}_q(\mathbb F_2).             \tag{F.IV4}
\]
No finite-index sublattice has been substituted for a norm
image; the displayed block calculations are surjective
integer formulas.

The centralizer action on (F.IV4) is precisely independent
similarity in the two matrix factors.  Indeed, in a product
of commuting block matrices, every contribution factoring
through a swap block is a norm.  The remaining diagonal
blocks are just the reductions of the trivial and sign
blocks of the centralizer matrix.  If that matrix is an
integral unit, its inverse shows that both reductions are
invertible.  Conversely block-diagonal unimodular changes
on the trivial and sign summands induce every element of
$\mathrm{GL}_p(\mathbb F_2)\times
\mathrm{GL}_q(\mathbb F_2)$; elementary matrices give
unimodular lifts.  There is no additional identification
between the two factors, since an intertwiner of a trivial
and a sign lattice is zero over $\mathbb Z$.

It remains a genuinely finite classification.  Enumerate
all $2^{n^2}$ matrices of size $n=0,1,2,3$ over $\mathbb F_2$
and conjugate by every invertible matrix of that size.
Their respective numbers of similarity classes are
$1,2,6,14$.  Taking the smallest bit-coded representative
in each orbit gives a fixed canonical choice.
The six integral involution types and their counts are
\[
\begin{array}{c|rrrrrr}
 (p,q,r)&(0,3,0)&(1,2,0)&(2,1,0)&(3,0,0)&(0,1,1)&(1,0,1)\\\hline
 \text{classes}&14&12&12&14&2&2
\end{array}
\]
and their sum is 56.  A representative is (F.IV2) with
$B$ equal to the chosen $0$--$1$ representatives in its
trivial and sign blocks, and zero in every other block.

This also gives an effective normalizer for an arbitrary
input, not merely a list of representatives.  Normalize its
cusp and then its involution by the preceding Smith and
glue constructions.  Use the finite-field orbit lookup to
normalize its two reductions.  In dimensions at most three
the binary lift of an invertible binary matrix is already
unimodular: its determinant is odd and has absolute value
at most two.  The latter bound follows by adjoining a row
and column of ones and applying the $4\times4$ Hadamard
bound.  Finally solve the full integral linear equation
$AD+DA=B-B_{\rm can}$ by Smith reduction.  The preceding
norm calculation proves solvability.  Conjugating by
$\begin{pmatrix}I&D\\0&I\end{pmatrix}$ subtracts this norm
from the upper-right block.  The product of these explicitly
recorded conjugators is integral unimodular and carries all
three marked matrices to their canonical representative.

Both matrices in (F.IV2) have a three-dimensional fixed
space.  Choose a primitive integral fixed vector $v_i$.
Then $g_i(x)=M_ix+v_i/4$ has actual order four:
$g_i^2$ is the nontrivial translation $v_i/2$.
This order-two subgroup is free, and every nontrivial
subgroup of a cyclic group of order four contains it, so
the full action is free.  For any other finite affine lift
the actual order is still even, since its linear part has
order two.  The two-point orbifold quotient of F.2.31 is
therefore the nontrivial cyclic group of order
$\gcd(m_1,m_2)$.  Simple connectivity is impossible for
every choice of twists and actual orders.
\end{proof}

The complete finite enumeration and arbitrary-input
normalizer are implemented in
\path{code/fourfolds_involution_classes.py}.
The program verifies the full norm-quotient Smith form in
each of the six integral types, all finite-field similarity
orbits, all 56 representatives, and a nontrivial global
integral conjugate of each representative.  It returns and
checks 112 unimodular conjugators.  It also checks the 112
chosen order-four affine lifts and computes the complete
fundamental-group structure and $H_1$ of all 56 fillings
using F.2.31.  The exact matrices, numerators, orbit tables
and group data are saved in
\path{records/fourfolds_involution_classes.json}.
These executed finite sets exhaust this involution stratum;
they are not a height truncation of the other finite-order
strata.
